% preamble.tex
\usepackage[T2A]{fontenc}
\usepackage[utf8]{inputenc}
\usepackage[russian]{babel}
\usepackage{amsmath}
\usepackage{amssymb}
\usepackage{amsthm}
\usepackage{geometry}
\usepackage{hyperref}
\usepackage{graphicx}
\usepackage{float}

\geometry{a4paper, margin=1in}

% Отключаем курсив по умолчанию для всех теоремных окружений
\theoremstyle{definition}
\newtheorem{definition}{Определение}[section]
\newtheorem{theorem}{Теорема}[section]
\newtheorem{property}{Свойство}[section]
\newtheorem{corollary}{Следствие}[section]
\newtheorem{example}{Пример}[section]
\newtheorem{remark}{Ремарка}[section]
\newtheorem{lemma}{Лемма}[section]

% Настройка гиперссылок и TOC
\hypersetup{
    colorlinks=true,
    linkcolor=blue,
    citecolor=blue,
    urlcolor=blue,
    pdftitle={Билеты по физике},
    pdfauthor={Студент}
}
